\chapter{Proofs for dynamic-game completeness and its correction}
\label{annex:p18-is-dynamic-completeness}

\begin{annexitem}{annex:p18.is-compression-countermodel}{thm:p18-is-compression-countermodel}{All-arguments calculation of the finite countermodel}
\begin{proof}
First record the complete possible forms of strict consequences.
No strict rule has an empty premise, so \(C_S([])=\varnothing\).
All positive ordinary atoms come from \(u,v,t,z\), respectively.
The only positive priority conclusions supplied by defeasible rules
are \(m\triangleleft t\) and \(n\triangleleft z\).
Their four endpoints are distinct. Transitivity cannot compose
these two edges, and each of the other strict priority rules has
a negative conclusion. Induction on strict derivations therefore
shows that these are the only possible positive priority conclusions
in any argument or its strict extensions. The first requires \(p\)
and hence \(a\); the second requires \(q\) and hence \(b\).
The negative priority consequences are only the reverse negative
edges. In particular, no argument produces
\(\neg(m\triangleleft t)\) or \(\neg(n\triangleleft z)\).

The ordinary negative atom \(\neg a\) has only the strict producer
with premises \(a,b\). Any argument whose strict closure contains
\(\neg a\) therefore has both \(a,b\) and is incoherent.
The atom \(\neg c\) is produced either directly by \(m\) or by
the strict rule with premises \(c,a\). In the latter case the
argument is incoherent, since it already contains \(c\).
Likewise \(\neg d\) is produced directly by \(n\) or through the
incoherent combination \(d,b\). These alternatives exhaust the
ordinary strict conclusions. There are no weak premises, so
undercutting never occurs in this countertheory.

The empty argument defeats every incoherent argument, strictly:
none can defeat the empty argument, which is coherent and has
no conclusion or assumption. It is itself unattacked and is in
\(J_1\). Neither \(V\) nor \(Q\) has any attacker: there is no
producer of \(\neg b\) or of a contrary to its priority consequences.
They too are in \(J_1\).
Any attacker of \(U\) or \(P\) must produce \(\neg a\);
no contrary priority conclusion is available. Every such attacker
is incoherent by the preceding producer analysis and is strictly
defeated by the already justified empty argument.
Consequently \(U,P\in J_2\).

Consider \(Z\) relative to \(J_1\).
An attacker concluding \(\neg d\) directly with \(n\) has frontier
\(\{n\}\), which is lower than \(\{z\}\) because \(Q\in J_1\).
That rebuttal is suppressed. Any remaining rebuttal must derive
\(\neg d\) through the strict \(d,b\) rule and is incoherent.
The empty argument strictly defends against it.
This proves \(Z\in J_2\), accounting for all strict extensions
and for contents having more than one relevant conclusion.

At support \(J_2\) both positive priorities are established.
For \(T\) and \(A\), every direct \(\neg c\) or \(\neg d\)
rebuttal is consequently suppressed by its frontier comparison.
All remaining contrary conclusions use the strict producers
already shown to be incoherent. Each actual remaining defeater
is strictly defeated by the empty argument.
Both \(T\) and \(A\) are coherent themselves: their strict closures
contain no \(a,b\) and no ordinary contradiction.
Thus \(T,A\in J_3\). The inclusions asserted in the theorem
do not omit combined or strictly extended arguments from the
universe; the classification of every possible contrary producer
has covered them.

The strict closures of \(P\) and \(Q\), separately, are compatible:
their ordinary positive conclusions are \(a,b\), respectively,
and their two disjoint priority edges have no contrary between
the two closures. Neither contains an assumption. Thus
\(\{P,Q\}\) is conflict-free.
Any single argument containing both positive priority conclusions,
however, has both \(a,b\), hence a strict extension concluding
\(\neg a\). It is incoherent. The fact that the arguments are
pairwise compatible does not license that combined content.

The argument \(B=[m,n]\) has strict closure
\(\{\neg c,\neg d\}\), so is coherent.
With empty comparison it rebuts \(A\) through either \(c\) or \(d\).
Hence \(B D_\varnothing A\).
Suppose a response \(C\) neutralizes this challenge.
Failure of even the particular \(\neg c,c\) rebuttal witness
requires
\(\{m\}<_{\{C\}}\{t\}\), and failure of the \(\neg d,d\)
witness requires
\(\{n\}<_{\{C\}}\{z\}\).
Thus \(C\) concludes both positive priorities. The preceding
paragraph proves it incoherent.

Alternatively suppose \(C D_{\{C\}}^s B\).
The empty incoherence clause cannot give a defeat of the coherent
\(B\), and there are no undercuts.
Therefore \(C\)'s strict closure must contain \(c\) or \(d\),
the only possible complementary conclusions to those of \(B\).
If it contains \(c\), its producer is \(t\), with frontier
\(\{t\}\). Since \(B\) must not defeat \(C\), the particular
reverse witness with \(m\) must be suppressed.
Therefore \(C\) concludes \(m\triangleleft t\), which requires
the ordinary conclusion \(a\).
The strict rule \(c,a\to\neg c\) makes \(C\) incoherent.
If its forward witness uses \(d\), the same argument forces
\(n\triangleleft z\), then \(b\), and the strict contradiction
from \(d,b\to\neg d\).
Every possible response has therefore been covered.

After any such response the opponent legally plays \([]\).
No argument strictly defeats \([]\), for its strict closure and
assumption set are empty and it is coherent.
Nor can a priority response neutralize the designated defeat
of the incoherent preceding \(C\): that clause is independent
of priority support.
Thus the proponent has no legal next response.
The opponent forces loss whether or not either global minimality
or proponent nonrepetition is imposed.

If \(u':\Rightarrow a\) is added, the only change in this reasoning
is that \(a\) can have either of two producers. Every argument
containing \(p\) still contains \(a\), and every derivation of
\(\neg a\) still requires both \(a,b\).
All incoherence and response arguments therefore persist.
Each legal argument chooses a unique producer for its head;
an invalid content containing both \(u,u'\) is never needed
or counted as a response.
\end{proof}
\end{annexitem}

\begin{annexitem}{annex:p18.is-compression-rank}{thm:p18-is-compression-rank}{The precise failed merger inference}
\begin{proof}
The argument \(Z\) conflicts with \(B\) only at \(d,\neg d\).
At support \(\{Q\}\), \(n\triangleleft z\) holds and its reverse
does not. Thus \(Z\) rebuts \(B\), the reverse rebuttal is
suppressed, and \(Z D_{\{Q\}}^s B\).
The closure of \(Q\) contains \(b,n\triangleleft z\) and its
reverse negative, none contrary to \(\neg c,\neg d\).
Hence \(B\) does not attack \(Q\).

The list \([z,v,q]\) is legal: its distinct heads are \(d,b\)
and \(n\triangleleft z\), and \(b\) occurs before the rule
requiring it. Its own positive priority suppresses \(B\)'s
only reverse conflict, at \(d\), while its \(d\) conclusion
rebuts \(B\). Therefore
\((Z+Q)D_{\{Z+Q\}}^s B\).
But its strict closure contains \(\neg d\), so it is incoherent.
An incoherent argument cannot enter \(J\): the justified empty
argument defeats it, and no argument can strictly defeat that
empty defeater. This proves \(Z+Q\notin J\).

The previous theorem gives \(Q\in J_1\), \(Z\in J_2\).
In fact \(Z\notin J_1\), since the coherent \([n]\) defeats it
at empty comparison. No member of \(J_1\) can conflict with \(B\):
such a member would have to conclude \(c\) or \(d\), and every
such argument has an empty-comparison defeater \([m]\) or
\([n]\), unless already incoherent and defeated by \([]\).
There are no weak antecedents that could block those defeats.
Thus even selecting the first stage at which an argument
conflicts with \(B\) does not repair the merger inference.
The selected \(Z\) and its earlier support \(Q\) have exactly
that stage relation, yet their legal merger is never justified.
\end{proof}
\end{annexitem}

\begin{annexitem}{annex:p18.is-licensing-sound}{thm:p18-is-licensing-sound}{An opponent strategy against every unjustified root}
\begin{proof}
All reasoning is on the original dynamic justified set \(J\).
For each \(C\notin J=G(J)\), choose an argument \(B_C\) satisfying
\[
 B_C D_J C,\qquad
 \text{there is no }H\in J\text{ with }H D_J^s B_C.
\]
Such a witness exists by the definition of \(G(J)\).
For set-sized theories these simultaneous choices use ordinary
set-theoretic choice. They constitute a mathematical strategy,
not an assertion that \(J\) is effectively computable.

The opponent acts as follows. At a hearing whose claim \(C\)
is outside \(J\), choose \(B_C\).
It is legal because \(D_J\subseteq D_\varnothing\).
At a claim in \(J\), choose any legal opponent challenge if
one exists. If none exists the automatic closing rule applies.
We show that along this strategy every licensed entry is in \(J\).

Initially the list is empty. Suppose, toward the first contrary
admission, that the current list satisfies \(|L|\subseteq J\).
A claim \(C\notin J\) cannot close automatically: the selected
\(B_C\) is an empty-comparison defeater.
It cannot close by a defense \(H\in |L|\), since strict stability
would take \(H D_L^sB_C\) to \(H D_J^sB_C\), contrary to the
choice of \(B_C\).
It cannot close by neutralization, since
\(B_CD_JC\) and \(D_J\subseteq D_L\) rule that out.
Opening further hearings changes no licensed entry until some
frame closes. The same argument applies to the first purported
closure outside \(J\), wherever that frame lies in the stack.
Every closure therefore preserves \(|L|\subseteq J\).

If the root \(A_0\notin J\), it can never close under this
opponent strategy. A finite maximal play then ends with no
proponent move and is a loss; an infinite play is a loss by
the declared convention. Thus the opponent wins, and the
existence of a proponent winning strategy implies \(A_0\in J\).
The argument deliberately establishes the licensed-list invariant
for the selected opponent strategy. A record produced under
arbitrary cooperative opponent choices is not treated as an
independently verified collection of justified lemmas.
\end{proof}
\end{annexitem}

\begin{annexitem}{annex:p18.is-licensing-complete}{thm:p18-is-licensing-complete}{A terminating strategy from semantic entry ranks}
\begin{proof}
For \(C\in J\), let \(\rho(C)\) be the least ordinal \(\beta\)
with \(C\in J_\beta\). It is a nonzero successor ordinal:
an element of a limit union already occurs earlier.
We prove the following stronger assertion by transfinite induction
on \(\rho(C)\). Suppose \(C\) is the active claim, the entry set
of the licensed list is a subset of \(J\), and every suspended ancestor claim
has strictly larger entry rank. Then the proponent has a strategy
that closes the hearing for \(C\) after finitely many moves
on every play, licenses only members of \(J\), and leaves all
ancestor frames available for their prescribed return.

If the opponent has no \(D_\varnothing\)-challenge, the claim
closes automatically. Otherwise let its chosen challenge be \(B\).
Write \(\rho(C)=\alpha+1\). Since \(C\in G(J_\alpha)\), either
\(B\) no longer defeats \(C\) at \(J_\alpha\), or it has some
strict defender \(H_0\in J_\alpha\).

In the first case, defeat between the fixed finite arguments
\(B,C\) depends on finitely many priority bits, by
Theorem~\ref{thm:p18-is-dynamic-finitary}.
For every bit present in \(E_{J_\alpha}\) among those bits,
choose one witnessing argument from \(J_\alpha\).
Let \(X\) be their finite set. The same bits are present in
\(E_X\); absent bits cannot be introduced because
\(X\subseteq J_\alpha\). Hence \(\neg(B D_X C)\).
In particular \(X\ne\varnothing\): otherwise this would
contradict the actual empty-comparison challenge.

In the second case, apply the same finite-information argument
to both defeat directions between \(H_0\) and \(B\).
Choose their finitely many present-bit witnesses in \(J_\alpha\),
and include \(H_0\). Their finite set \(X\subseteq J_\alpha\)
satisfies \(H_0 D_X^s B\).
This is a finite choice of separate argument witnesses. No
assumption about a legal or coherent union has been made.

The proponent now opens hearings, one at a time, for the members
of \(X\) not already licensed. Each such member \(H\) has
\(\rho(H)<\rho(C)\). It is not an ancestor claim, all of whose
ranks are larger than that of \(C\), and it is not the current
claim. Its hearing is therefore fresh and legal.
By the induction hypothesis its hearing terminates on every
play and licenses only justified arguments. Upon its closure,
the exact parent challenge \((C,B)\) resumes.
There are finitely many members of \(X\), so finitely many such
terminating hearings suffice. Already licensed members are
simply skipped, even if their arguments share heads or
subarguments with another member.

At the end let \(L'\) be the licensed list. Its entry set contains
\(X\) and is a subset of conflict-free \(J\).
In the neutralization case, antitonicity transfers
\(\neg(B D_XC)\) to \(\neg(B D_{L'}C)\); any member of
nonempty \(X\) is a legal citation.
In the defense case, strict stability transfers
\(H_0D_X^sB\) to \(H_0D_{L'}^sB\).
The appropriate resolution closes \(C\), which is itself in \(J\),
and returns to its parent. This proves the induction assertion.

The induction also proves termination with its full quantifiers:
for every selected \(B\), there are finitely many proposed
hearings, and each has a strategy terminating on every possible
opponent continuation. A finite succession of finite completed
plays is finite. No uniform finite bound across different
opponent choices is required in the general theorem.
Nested fresh hearings strictly decrease ordinal entry rank;
no unresolved claim can be used to establish itself.

Apply the assertion to the initial empty list and the root
\(A_0\in J\), with no ancestors. It supplies a winning strategy.
For roots outside \(J\), the previous proof supplies an opponent
winning strategy. These conclusions establish the equivalence
and the specified determinacy without assuming a merger property,
global minimality, or finite semantic rank.
If \(A'\) is a source subargument of a winning claim \(A_0\),
soundness gives \(A_0\in J\), subargument closure of the dynamic
semantics gives \(A'\in J\), and the completeness construction
provides its winning strategy. This proves the stated closure
property without an independent syntactic merger claim.
\end{proof}
\end{annexitem}

\begin{annexitem}{annex:p18.is-licensing-finite}{thm:p18-is-licensing-finite}{Finite strategy trees and finite-domain computation}
\begin{proof}
Fix the winning proponent strategy constructed above.
At proponent positions it has one selected move; the automatic
closing and return transitions also have one successor.
At an opponent position the possible moves are the
\(D_\varnothing\)-defeaters of a fixed argument.
In a finitary theory that is a finite set: defeat implies attack
except for the possible additional empty incoherence defeater.
Thus the strategy tree is finitely branching.

It has no infinite branch by the completeness proof.
A finitely branching rooted tree with infinitely many nodes
would have an infinite branch: from a node with infinitely
many descendants, one of its finitely many children has
infinitely many descendants; choose it and repeat.
Starting at the root yields a contradictory infinite path.
Hence the winning strategy tree is finite.
The reverse implication follows from general soundness.

If the universe has \(N<\infty\) arguments, every opened claim
is distinct from all pending and all previously licensed claims.
A claim is never removed except by closing, which licenses it.
Therefore at most \(N\) hearings, counting the initial root,
can ever be opened during one play.
Each hearing receives at most one opponent challenge, opens
children only through those fresh-hearing moves, and has at most
one closing response. The entire play consequently has a finite
bound in terms of \(N\), even for poorly chosen strategies.
There is no way to obtain infinitely many moves merely by
repeated citation: a citation closes its active frame.

For finitely many ground rules, the prior finite-content
construction enumerates the legal universe and computes every
strict closure, frontier and \(D_L\) comparison.
The licensed list and pending stack are sequences without
repetition drawn from that finite universe; each frame stores
at most one opponent argument. The state space is finite,
and the preceding move bound permits backward induction from
the terminal outcomes. This computes the winner and choices
for a winning strategy. It uses actual finite grounding,
not an unsupported inference from finite rule schemata or
from abstract finitariness.
\end{proof}
\end{annexitem}

\begin{annexitem}{annex:p18.is-licensing-worked}{thm:p18-is-licensing-worked}{The complete repaired counterexample dialogue}
\begin{proof}
The initial hearings for \([]\) and \(Q\) close automatically
because they have no empty-comparison defeaters.
Every challenge to \(P=[u,p]\) must produce \(\neg a\);
no contrary priority producer exists.
The countermodel proof showed that each such argument is
incoherent and is strictly defeated by \([]\), independently
of priority context. Thus every opponent choice in this hearing
is met by the same licensed defense. The hearing closes and
licenses \(P\).

Now \(E_L\) contains \(m\triangleleft t\) and
\(n\triangleleft z\). For the particular initial
\(B=[m,n]\), both contrary frontier comparisons are suppressed,
so \(\neg(BD_LA)\). A neutralization citation of \(P\) is legal.
It closes the root, rather than merely declaring the earlier
lemma hearing successful.

For an arbitrary initial opponent \(B'\), any conflict with
\(A\) is at \(c\) or \(d\).
If a surviving contrary derivation uses one of the ordinary
strict rules, \(B'\) is incoherent and the proponent cites
the licensed empty argument as strict defender.
Otherwise all its contrary derivations use the direct
defaults \(m,n\); both comparisons are suppressed after
licensing \(P,Q\), and the proponent neutralizes.
If an argument has both kinds of witness, the incoherence
case already supplies the defense.
No other attack form is possible because there are no weak
premises, no contrary to a different root conclusion, and
the root itself is coherent. This gives all branches of
the winning strategy.

The optional hearing for \(U\) has the same incoherent
challengers as the ordinary \(a\) part of \(P\); it closes
by citing the already licensed empty argument.
The later hearing for \(P\) is still fresh: \(P\ne U\),
although \(U\) is its subargument.
No rule forbids this overlap, and no union is formed.
With the alternative producer \(u'\), choose either legal
argument \([u,p]\) or \([u',p]\). Each hearing has exactly
the same exhaustive kind of challenge and resolution.
Distinct licensed entries may therefore share a conclusion
without violating the unique-producer condition inside each
individual argument.

For the separate cooperative-play example, the empty argument
and \([r_b]\) are unattacked. Every coherent argument containing
\(r_s\) is strictly undercut by \([r_b]\); a content containing
both \(r_s\) and \(r_b\) is incoherent.
Every coherent argument concluding \(a\) is rebuffed by the
coherent singleton \([r_{\neg a}]\), which neither the empty
argument nor \([r_b]\) defeats.
Every coherent argument concluding \(\neg a\) through
\(r_{\neg a}\) likewise has a coherent \(a\)-attacker.
No other positive producer is available; combining both
ordinary contrary conclusions makes the content incoherent.
These observations show directly that the first-stage set
\(\{[],[r_b]\}\) is already a fixed point and is \(J\).
The response to the selected challenge \([r_s]\) is nevertheless
legal: the \([r_b]\) hearing closes automatically, and its
undercut of \([r_s]\) is strict.
This proves both the completed-play calculation and the absence
of a winning strategy for the unjustified root, the latter also
following from the general opponent construction.
\end{proof}
\end{annexitem}

\begin{annexitem}{annex:p18.is-licensing-reversed}{thm:p18-is-licensing-reversed}{Dynamic overruling and residual defensibility}
\begin{proof}
If \(B\in J\) attacks \(A_0\), the challenger chooses that
argument and then uses its winning licensing strategy.
This wins the challenge protocol against every defender response.
Conversely any challenger winning strategy must first choose
an attacker \(B\), and its continuation must win the licensing
dialogue for \(B\). Soundness puts \(B\) in \(J\).
This proves the exact equivalence with source overruling.

If no justified argument attacks \(A_0\), every permitted initial
witness \(B\) lies outside \(J\). The defender then uses the
opponent strategy constructed in
Theorem~\ref{thm:p18-is-licensing-sound} for that chosen \(B\).
It forces a finite refusal or an infinite unsuccessful proving
play, each a defender win. If there is no initial attacker,
the challenger loses without a hearing.
Thus the outcome is determined also in this case.

By definition an argument is source-defensible when it is
neither justified nor overruled. Substituting the just-proved
characterization gives the stated conjunction.
If \(A_0\in J\), conflict-freedom of \(J\) excludes any justified
attacker, so the defender wins as well; no stronger residual
classification follows from that outcome alone.

In a finitary theory a chosen justified attack witness has a
finite licensing proof by the preceding theorem.
For a defender win, however, the proving player may successively
open hearings for unrelated arguments on an infinite universe.
The declared loss of infinite proving play makes this a defender
win without requiring a finite refutation tree.
No finite negative certificate or uniform decision algorithm
is inferred. On a finite ground-rule domain the attack choices
and the subsequent licensing game are finite and decidable.

Finally the root \(A\) of the countertheory belongs to \(J\),
and its \(c,d\) conclusions attack \(B=[m,n]\).
These facts prove the two worked status claims.
For the earlier incoherent two-default combination, the empty
argument has no conclusion or assumption and hence is not
an attack witness. The designated incoherence defeat does
not change the protocol's stated attack requirement.
\end{proof}
\end{annexitem}
